\documentclass[french]{book}
\usepackage{teach}
\usepackage{verbatim}
\usepackage{amsmath,amsfonts,amssymb}
\usepackage{pgfplots}
%\usepackage{geometry}
%\geometry{top=1cm, bottom=1cm, left=2cm, right=2cm}
\usepackage[utf8]{inputenc}
\usepackage[french]{babel}
\redefinecolor{thm}{title}{bgcolor}{alizarin}
\redefinecolor{prop}{title}{bgcolor}{meatbrown}
\redefinecolor{defn}{title}{bgcolor}{olivine}
\definecolor{alizarin}{rgb}{0.82, 0.1, 0.26}
\definecolor{meatbrown}{rgb}{0.9, 0.72, 0.23}
\definecolor{olivine}{rgb}{0.6, 0.73, 0.45}
\usepackage{pgf,tikz,tkz-tab}

\usepackage{mathrsfs}
\usetikzlibrary{arrows}

\begin{document}
\pagestyle{empty}
\section*{Devoir à la maison}
L'objectif de ce devoir est de comprendre et d'observer l'impact qu'on les opérations appliquées à une fonction sur une représentation graphique. Tout le long de ce devoir il est demandé d'être sur le site desmos.com pour représenter les fonctions plus rapidement. 
\subsection*{Niveau 1 - Fonction affine et droite}

\begin{defn}
Une fonction affine est une fonction $f$ définie par $f(x)=ax+b$, où $a,b\in\mathbb{R}$.
\vspace{5mm}
\end{defn}

\subsubsection*{Exemple}
Soit $f$ la fonction définie par $f(x)=-2x+3$
\begin{enumerate}
\item Déterminer $D_f$, faire un tableau de valeurs avec au moins 5 valeurs distinctes et donner la représentation graphique de $f$ dans un repère orthonormée
\item Dresser le tableau de variation et le tableau de signe de $f$.
\end{enumerate}

\subsubsection*{Conjecture}
\begin{enumerate}
\item Rentrer dans desmos $f(x)=ax+b$ en ajoutant les curseurs $a,b$ sur desmos puis faite varier $a$ et $b$. 
\item Que peut-on dire des variations de $f$ en fonction de $a$? et de $b$?
\end{enumerate}

\subsubsection*{Preuve}
\begin{enumerate}
\item Soit $x<y$ deux réels, montrer que si $a>0$ alors $f(x)<f(y)$ (c'est à dire que $f$ est strictement croissante sur $\mathbb{R}$).
\item Soit $x<y$ deux réels, montrer que si $a<0$ alors $f(x)>f(y)$ (c'est à dire que $f$ est strictement décroissante sur $\mathbb{R}$).
\end{enumerate}

\subsection*{Niveau 2 - Fonction du second degré et parabole}


\begin{defn}
Une fonction du second degré est une fonction définie par  $f(x)=ax^2+bx+c$, où $a\in \mathbb{R}^*$ et $b,c\in\mathbb{R}$.
\vspace{5mm}
\end{defn}


\subsubsection*{Exemple}
Soit $f$ la fonction définie par $f(x)=(-2x+3)^2-1$
\begin{enumerate}
\item Montrer que $f$ est une fonction du second degré.
\item Déterminer $D_f$, faire un tableau de valeurs avec au moins 5 valeurs distinctes et donner la représentation graphique de $f$ dans un repère orthonormée.
\item Résoudre algébriquement $f(x)=0$
\item Dresser le tableau de variation et le tableau de signe de $f$ à l'aide de la représentation graphique de $f$ sur desmos.
\end{enumerate}

\subsubsection*{Conjecture}
\begin{enumerate}
\item Rentrer dans desmos $f(x)=a(x-b)^2+c$ en ajoutant les curseurs $a,b,c$ sur desmos puis faite varier $a,b$ et $c$. 
\item Soit $a>0$, que peut-on dire du nombre de solution à l'équation $f(x)=0$ en fonction de $c$?
\item Quelle influence à $b$ sur le placement de la courbe représentative de $f$? Même question pour $c$.

\end{enumerate}

\subsubsection*{Preuve}
\begin{enumerate}
\item Soit $a>0$ et $c<0$ exprimer les solutions de $f(x)=0$ en fonction de $a,b$ et $c$. 
\item Soit $g(x)=x^2-7x+12$ déterminer $a,b$ et $c$ tel que $f(x)=g(x)$, puis résoudre $g(x)=0$ 
\end{enumerate}

\subsection*{Niveau 3 - Fonction homographique et hyperbole}


\begin{defn}
Une fonction homographique est une fonction définie par  $f(x)=\dfrac{ax+b}{cx+d}$, où $a,b,c,d\in\mathbb{R}$ et tel que $ad-bc \neq 0$.
\vspace{5mm}
\end{defn}


\subsubsection*{Exemple}
Soit $f$ la fonction définie par $f(x)=\dfrac{-2x+3}{x-4}$
\begin{enumerate}
\item Déterminer $D_f$, faire un tableau de valeurs avec au moins 5 valeurs distinctes et donner la représentation graphique de $f$ dans un repère orthonormée.
\item Résoudre algébriquement $f(x)=0$
\item Dresser le tableau de variation à l'aide de la représentation graphique de $f$ sur desmos et le tableau de signe sans l'aide de desmos.
\end{enumerate}

\subsubsection*{Conjecture}
\begin{enumerate}
\item Rentrer dans desmos $f(x)=\dfrac{a}{x-b}+c$ en ajoutant les curseurs $a,b,c$ sur desmos puis faite varier $a,b$ et $c$. 
\item Soit $a>0$, que peut-on dire des variation de $f$?
\item Quelle influence à $b$ sur le placement de la courbe représentative de $f$? Même question pour $c$.

\end{enumerate}

\subsubsection*{Preuve}
\begin{enumerate}
\item Montrer que si $a>0$, alors $f$ est strictement décroissante.
\item Soit $g(x)=\dfrac{-2x+3}{x-4}$ déterminer $a,b$ et $c$ tel que $f(x)=g(x)$.
\end{enumerate}

\end{document}